\chapter{时钟、复位与同步边界}

如果每个存储单元随时都能改变，大电路很快会混乱。同步电路采用统一节奏：一个时钟沿让一批状态更新，随后组合逻辑根据这些状态传播，直到下一个时钟沿再保存新结果。时钟不是装饰，也不只是“频率数字”，它是状态更新边界。

\section{一、时钟周期把计算和保存分开}

先看一条最基本的寄存器到寄存器路径：

\begin{lstlisting}
source register -> combinational logic -> destination register
\end{lstlisting}

时钟沿到来后，源寄存器输出旧状态对应的新值。这个值进入组合逻辑，经过门延迟和布线延迟，最终到达目标寄存器输入。下一个时钟沿到来时，目标寄存器采样这个结果。

\begin{pdfdiagram}[x=1cm,y=1cm]
  \node[hw state,minimum width=2.0cm] (src) at (1.2,1.35) {源寄存器};
  \node[hw compute,minimum width=2.35cm] (logic) at (4.55,1.35) {组合逻辑};
  \node[hw state,minimum width=2.0cm] (dst) at (8.1,1.35) {目标寄存器};
  \draw[hw bus] (src.east) -- node[above,hw label] {clock-to-Q} (logic.west);
  \draw[hw bus] (logic.east) -- node[above,hw label] {逻辑 + 布线延迟} (dst.west);
  \node[hw label,align=center] at (1.2,0.35) {时钟沿后\\输出新状态};
  \node[hw label,align=center] at (4.55,0.35) {一个周期内\\传播并稳定};
  \node[hw label,align=center] at (8.1,0.35) {下个时钟沿\\采样结果};
\end{pdfdiagram}
\diagramnote{同步路径的周期预算来自寄存器输出延迟、组合逻辑延迟、布线延迟和目标寄存器建立时间。}

\begin{lstlisting}
after clock edge:
  source register output changes
between edges:
  combinational logic settles
at next edge:
  destination register captures
\end{lstlisting}

因此频率不是越高越好。频率提高意味着周期缩短，留给组合逻辑稳定的时间变少。若周期短到最慢路径来不及稳定，目标寄存器就可能采到旧值、过渡值或错误值。

\section{二、建立时间和保持时间是采样窗口}

触发器采样不是数学瞬间。输入 D 必须在时钟沿到来前稳定一小段时间，这叫建立时间；时钟沿之后还要继续稳定一小段时间，这叫保持时间。

\begin{longtable}{p{0.22\linewidth}p{0.42\linewidth}p{0.26\linewidth}}
\toprule
约束 & 含义 & 违反后果 \\
\midrule
建立时间 & 采样前输入要提前稳定 & 采到错误值或不稳定值 \\
保持时间 & 采样后输入要继续稳定 & 新旧值混入采样窗口 \\
传播延迟 & 上一级输出和组合逻辑需要时间 & 周期太短会来不及 \\
\bottomrule
\end{longtable}

举一个具体路径。源寄存器在时钟沿后并不会立刻输出新值，它有一个 clock-to-Q 延迟；组合逻辑本身有延迟；线也有延迟。目标寄存器又要求输入在下一个时钟沿前提前稳定。时钟周期必须覆盖这些时间。

\begin{lstlisting}
clock period >= source clock-to-Q
              + combinational delay
              + wire delay
              + destination setup time
\end{lstlisting}

保持时间看起来更小，却同样危险。若源寄存器输出太快、路径太短，目标寄存器可能在同一个时钟沿之后立刻看到新值，从而破坏它本应采旧值的窗口。长路径会限制最高频率，短路径也可能破坏保持约束。

\section{三、关键路径决定最高可用频率}

一块同步电路里会有很多寄存器到寄存器路径。最长、最慢、最限制周期的那条叫关键路径。它可能是一串加法器，也可能是多级选择器，也可能是远距离布线加上复杂译码。

例如一个 32-bit 串行进位加法器接在两个寄存器之间。最低位进位要一位一位传到最高位，最高位结果最后才稳定。如果把这个加法器放在一个周期里完成，那么周期至少要等最坏进位链走完。若把运算拆成更短的几段，中间插入寄存器，单个周期可以变短，但完成一次完整运算需要更多周期。

\begin{pdfdiagram}[x=1cm,y=1cm]
  \node[hw label,text=BookNavy] at (2.8,2.5) {一周期完成：路径长};
  \node[hw state,minimum width=1.4cm] (r1) at (0.65,1.55) {R};
  \node[hw compute,minimum width=3.2cm] (addlong) at (3.0,1.55) {32-bit 加法器};
  \node[hw state,minimum width=1.4cm] (r2) at (5.45,1.55) {R};
  \draw[hw bus] (r1.east) -- (addlong.west);
  \draw[hw bus] (addlong.east) -- (r2.west);

  \node[hw label,text=BookNavy] at (8.5,2.5) {拆段完成：路径短但周期多};
  \node[hw state,minimum width=1.1cm] (p1) at (6.35,1.55) {R};
  \node[hw compute,minimum width=1.55cm] (a1) at (8.0,1.55) {低 16 位};
  \node[hw state,minimum width=1.1cm] (mid) at (9.65,1.55) {R};
  \node[hw compute,minimum width=1.55cm] (a2) at (11.3,1.55) {高 16 位};
  \node[hw state,minimum width=1.1cm] (p2) at (12.95,1.55) {R};
  \draw[hw bus] (p1.east) -- (a1.west);
  \draw[hw bus] (a1.east) -- (mid.west);
  \draw[hw bus] (mid.east) -- (a2.west);
  \draw[hw bus] (a2.east) -- (p2.west);
  \node[hw label,align=center] at (2.95,0.45) {周期必须覆盖整条加法路径};
  \node[hw label,align=center] at (9.65,0.45) {每段更短\\但需要跨更多边界};
\end{pdfdiagram}
\diagramnote{关键路径变短可以提高单周期频率，但插入状态边界会改变完成一次动作所需的周期数。}

这就是硬件设计常见取舍：更短周期通常要更多状态边界；更少周期通常要求更长组合路径。时钟不是单独设置出来的，它受电路结构约束。

\section{四、复位让状态从已知点开始}

触发器通电后的值不一定符合设计预期。若一组状态寄存器随机落在不同值，后面的组合逻辑可能立刻进入非法状态。复位信号的作用，是把关键状态单元带到已知点，例如计数器清零、状态机回到 idle、控制寄存器进入安全默认值。

复位本身也要讲边界。若复位释放时，有的触发器先退出复位，有的后退出复位，状态机可能看到一半新状态、一半旧状态。常见处理是异步拉低或拉高复位以便快速进入安全状态，再同步释放，让所有相关触发器在清楚的时钟边界恢复工作。

复位不是“清零所有东西”这么简单。哪些状态需要复位，哪些数据寄存器可以不复位，哪些输出在复位期间必须安全，都要由硬件需求决定。盲目复位太多节点会增加布线和时序压力；复位太少又可能让电路从未知状态开始。

\section{五、同步设计的基本纪律}

到这里，同步设计可以压成三条纪律。第一，状态只在明确边界更新。第二，组合逻辑必须在下一个边界前稳定。第三，跨边界信号必须先处理，不能假设外部变化刚好对齐本地时钟。

下一章会把这些纪律用于更大的状态结构。多个触发器并排形成寄存器；寄存器加反馈形成计数器；寄存器保存阶段编号，组合逻辑决定下一阶段，就形成有限状态机。

\begin{classicnote}
同步设计的核心标注是时序账本。每条从寄存器出发、经过组合逻辑、到达另一个寄存器的路径，都要交代传播延迟、建立时间、保持时间和时钟偏斜。

\begin{itemize}
\item 纸上实验：给一条路径写出 \code{Tclk >= tCQ + tcomb + tsetup + tskew}，再解释每一项对应哪个硬件现象。
\item 机器证据：关键路径决定最高频率，不是平均路径决定；一条慢路径就能限制整块同步电路。
\item 复位证据：复位要让控制状态进入已知安全点，但释放复位也要有同步边界，不能让状态机半边先跑。
\item 检查题：把一个 32-bit 加法器拆成两个 16-bit 阶段后，频率可能提高，但一条加法指令的总延迟一定降低吗？
\end{itemize}

经典教材强调“时间也是硬件资源”。面积、功能、时钟周期、吞吐和延迟互相牵制。一个逻辑表达式正确但错过采样窗口，在机器里仍然是错误电路。
\end{classicnote}
